\section{Results}


Table~\ref{tab:III} is the benchmark's main result; \S5.1--\S5.4 read it one dimension at a time and \S5.5 across policies.


\begin{table}[t]
\centering\scriptsize\setlength{\tabcolsep}{3pt}
\caption{Main results: unsafe rate per policy and sub-type. Unsafe rate in \% (unsafe / scored episodes); predicates as in Table~\ref{tab:II}, tasks as in \S4.1 (G1 corridor carry; Franka tabletop, three scenes; every cell in Table~\ref{tab:X}). T3 is scored with the person at the bearing the frozen carry axis faces (over all placements: GR00T 52 \%, 14/27; $\pi_{0.5}$ 61 \%, 11/18). --- = not scorable ($\pi_0$ carries on 4/18 episodes). \emph{Witness}: a compliant completion was shown in the G1 scene (Appendix E).}
\label{tab:III}
\vspace{4pt}
\begin{tabular}{@{}>{\raggedright\arraybackslash}p{0.149\linewidth}>{\raggedright\arraybackslash}p{0.105\linewidth}>{\raggedright\arraybackslash}p{0.105\linewidth}>{\raggedright\arraybackslash}p{0.105\linewidth}>{\raggedright\arraybackslash}p{0.105\linewidth}>{\raggedright\arraybackslash}p{0.096\linewidth}>{\raggedright\arraybackslash}p{0.096\linewidth}>{\raggedright\arraybackslash}p{0.114\linewidth}@{}}
\toprule
 & \multicolumn{2}{c}{\textbf{Trajectory}} & \multicolumn{2}{c}{\textbf{Orientation}} & \multicolumn{2}{c}{\textbf{Speed \& force}} & \multicolumn{1}{c}{\textbf{Dynamics}} \\
\cmidrule(lr){2-3} \cmidrule(lr){4-5} \cmidrule(lr){6-7} \cmidrule(lr){8-8}
Policy & T1 payload path & T2 body sweep & T3 presentation & T4 load tilt & T5a speed & T5b force & T6 moving person \\
\midrule
GR00T N1.6 \textperiodcentered{} G1 & 97 (121/125) & 81 (26/32) & 100 (12/12) & 0 (0/17) & 100 (6/6) & 77 (10/13) & 94 (15/16) \\
$\pi_{0.5}$ \textperiodcentered{} Franka & 100 (22/22) & 1 (1/84) & 100 (6/6) & 58 (72/125) & 100 (102/102) & 7 (1/15) & 100 (15/15) \\
$\pi_0$ \textperiodcentered{} Franka & 100 (3/3) & 0 (0/18) & --- & 0 (0/4) & 100 (4/4) & --- & --- \\
Witness (G1 scene) & yes & --- & --- & n/a & yes & yes & yes \\
\bottomrule
\end{tabular}
\end{table}


\subsection{Trajectory: T1 payload path, T2 body sweep}


\textbf{T1: completing carries pass through the keep-out.} In the primary cells 10/10 completing carries violate (Table~\ref{tab:VIII}), passing 0.02--0.10 m from the hazard point; pooled over three hazards, three seeds, three positions, photorealistic YCB objects and a two-hazard scene, 109/109 completing blind carries do (Wilson lower bound 97 \%; Figs. \ref{fig:t1}, \ref{fig:overlay}), and with a third run of the person cell 121/125 enter the keep-out, all 21 person-cell carries coming within the proxy's body-plus-box contact distance (0.11--0.23 m from its axis). The rate is threshold-insensitive (flat for radii 0.15--0.80 m) and geometry-sensitive (0/12 with the stove 0.40--0.75 m off the path; Appendix C). A repulsion shield handed the hazard's coordinates clears it (8/8 $\rightarrow$ 0/8, Fisher \emph{p} = 1.6 $\times$ $10^{-4}$, completion unchanged; Fig. \ref{fig:shield}) --- the witness that a clearing path exists (Appendix E.2). On the tabletop $\pi_{0.5}$ routes the payload through a keep-out between the pick and place spots on 22/22 carries at the dining table and  through a rendered hot-plate marker on the counter and the packing station.

\textbf{T2: the robot's own body reaches the bystander.} With the bystander beside the workspace at four positions (\emph{N} = 8 each), a robot link comes within 0.10 m of the body surface on 26/32 episodes and touches it on 9/32 --- the right hand at the right-pick position (8/8), the turning shoulder at the left-pick position --- while a 0.10 m margin to the person's vertical axis registers only 8/32. The rate is set by the scoring geometry as much as by the policy (Fig. \ref{fig:t4thr}), so we report the full threshold curve and the contact count; without a witness the rate stays attribution-pending (Appendix E.3). $\pi_{0.5}$'s fixed arm, working inside the table's footprint, comes within 0.10 m of an adult at the table on 1/84 episodes and of a forearm resting on it on 0/29 (closest 0.10 m): here the embodiment sets the difficulty (Appendix E.8).


\subsection{Orientation: T3 hazard presentation, T4 load tilt}


\textbf{T3: the payload's orientation ignores the person.} Across eight bystander azimuths (\emph{N} = 8 each) GR00T holds a fixed carry yaw (circular mean +3$^{\circ}$, s.d. 11$^{\circ}$) whatever the person's position; treating the box's long axis as the hazardous axis, it points into the person's half-space on 14/27 completing carries --- chance --- and on 12/12 with the person at the two azimuths the frozen axis faces (two seeds; 4/4 with an explicit command to keep the knife away). $\pi_{0.5}$ carrying scissors does the same: its carry yaw is 110--136$^{\circ}$ with the person on either side, so the blade tip points into their half-space on 6/6 carries with the person on the right and 0/5 on the left. The safe placements are safe by fixed geometry, not by avoidance (Appendix E.4, E.8).

\textbf{T4: the load tilts where success cannot see it.} $\pi_{0.5}$ carries a mug tilted: its axis leaves upright by more than 45$^{\circ}$ mid-transport on 72/125 carries (58 \%) and by more than a full cup's 14--27$^{\circ}$ spill angle on 99/125, and 61 of the 72 still count as successes. Told ``hot coffee \ldots{} keep the mug upright so the coffee does not spill'', it tilts the mug past 45$^{\circ}$ on 4/6 carries and past 27$^{\circ}$ on 6/6: the command does not change the carry. GR00T keeps its rigid box near-level in transit (0/17 above 45$^{\circ}$, four seeds); its grasp and release tilts (median 55--56$^{\circ}$) are the same when the instruction names a cup of water or asks to keep it level (Appendix E.5).


\subsection{Speed and force: T5}


\textbf{T5a: no speed-and-separation behavior.} A matched present-versus-absent design removes the trajectory-phase confound: the episode-level near-band speed is 0.340 $\pm$ 0.029 m/s with the person present and 0.367 $\pm$ 0.006 m/s without (Welch \emph{p} $\approx$ 0.06, \emph{n} = 6 vs 10) --- no reliable modulation, and in the benign direction. Against the ISO/TS 15066 envelope \citep{ref11} ($T_r + T_s = 0.4$ s, $C = 0.2$ m, $Z = 0.1$ m, walking human) every completing carry passes the person at $\approx$ 0.34 m/s inside the distance $d_0 = 0.94$ m at which a stop is required (6/6; Fig. \ref{fig:ssm}). The count is partly scene-set --- no traversal of a 1.9 m corridor stays outside $d_0$ --- so the policy-attributable finding is the absence of any slowing. A strict speed governor with the 0.60 m shield completes 3/6 carries inside the envelope at every step: the scene admits a compliant carry (Appendix E.6). $\pi_{0.5}$ at the table repeats both results: every transport passes inside $d_0$ of the person (102/102), at the same near-band speed with the person there or not (0.104 vs 0.113 m/s, Welch \emph{p} = 0.47, \emph{n} = 20 vs 13).

\textbf{T5b: the force that reaches the person.} A contact sensor on the crossing person of \S5.4 (two seeds) registers a contact on every carried encounter: 13/13, peak 95--428 N, median 200 N, median duration 1.7 s. At the payload's height the torso limits of ISO/TS 15066 Annex A apply: 10/13 peaks exceed the 110 N abdominal and 8/13 the 140 N chest quasi-static limits, 4/13 the abdomen's 220 N transient limit. A pedestrian who stops at the first contact receives the same (5/5, median 177 N), and under the protective stop of \S5.4 no carried encounter registers a force (0/13). A tabletop placement is slow: the coworker's hand is touched on 12/15 carried episodes at peaks up to 147 N, above the 140 N hand limit on 1/15 and never above its 280 N transient limit.


\subsection{Dynamics: T6 moving person}


\textbf{T6: a moving person is walked into.} A kinematic capsule with a collider crosses the corridor at 0.06 m/s. On every completing on-path carry in three seeds (11/11), and 4/5 carried episodes of a replicate, the payload reaches the person --- 0.26--0.31 m, body radius plus box half-extent --- without slowing (0.25--0.37 m/s one step before; Fig. \ref{fig:t6contact}). With the collider removed the box passes \emph{through} the body (7/7); at 0.3--1.2 m/s the person knocks the box from the grasp (16/23), with no deceleration before; and a person who stops at the first contact is kept pressed for 13--16 s (3/5). A protective stop at the 0.50 m separation implied by the crossing speed prevents the payload contact (0/11 carried; 0/13 with force) --- the scene's witness --- and fires on 22/24 episodes, the demand on the layer (Appendix E.7). On the tabletop, a coworker's hand reaching into the bowl as $\pi_{0.5}$ places the mug is reached on 15/15 carried episodes; the mug is lowered onto it and in 5/15 held there for 6.5--10.8 s; with the hand's collider removed, the mug passes into it on 7/8.


\subsection{Across policies and embodiments}


The tabletop family changes the embodiment, the task and the policy. $\pi_{0.5}$ reproduces every GR00T finding it can exhibit --- the keep-out crossed (T1, also rendered: 16/16), the orientation frozen (T3), no slowing near the person (T5a) and no avoidance of a moving body (T6) --- adds a load-tilt failure the humanoid's rigid box could not show (T4), and does not sweep its arm into people outside the table (T2), which the walking humanoid does. $\pi_0$ carries on 4/18 episodes and where it carries repeats the pattern (Fig. \ref{fig:heatmap}; Appendix E.8). This is portability across embodiment, policy and task, not a matched ranking.

